\lecture{19}
\title{Applications of the Spectral Theorem}

\begin{document}

Today will be dedicated to applying the Spectral Theorem in three broader contexts.

\subsection{Spectral Application \#1: Classifying Quadric Surfaces}

We've done enough foreshadowing already.
\begin{itemize}
    \item In Section 16.3, we foreshadowed how all quadratic polynomials $\mathbb{R}^n \to \mathbb{R}$ look like:
\[ \text{quadratic polynomials} = \underbrace{\left(\sum_{i, j} A_{i, j} x_i x_j\right)}_{\text{degree 2 terms}} + \underbrace{\left(\sum_{i} b_i x_i \right)}_{\text{degree 1 terms}} \ + \ \left(c\right) = x^{\top}Ax + b^{\top}x + c. \]
    \item In Section 17.5, we foreshadowed how the signature of $A$ encodes its behavior.
\end{itemize}
Let's get to the punchline of this foreshadowing.

\textbf{Definition (Quadric Surface).} A quadric surface in $\mathbb{R}^n$ is a subset of the form $\{x \in \mathbb{R}^n \mid x^{\top}Ax + b^{\top}x + c = 0 \}$ for some symmetric matrix $A \in \mathbb{R}^{n \times n}$, $b \in \mathbb{R}^n$, and $c \in \mathbb{R}$.

\textbf{Example.} We already know the $n = 1$ and $n = 2$ cases from high school.
\begin{itemize}
    \item In $n = 1$, quadric surfaces are either pairs of points, single points, the empty set, or all of $\mathbb{R}$.
    \item In $n = 2$, quadric surfaces are conics: ellipses, hyperbolas, and the degeneracies.
\end{itemize}
We now use a special case of the Spectral Theorem in the context of $\mathbb{R}$.

\begin{theorem}{Spectral Theorem over $\mathbb{R}^n$}
    If $A \in \mathbb{R}^{n \times n}$ is symmetric, its eigenvectors form an orthonormal basis.
\end{theorem}

In other words, there exists some $T \in O_n(\mathbb{R})$ such that $T^{\top}AT = \Lambda$, where $\Lambda$ is diagonal.  Thus,
\begin{align*}
    & x^{\top}Ax + b^{\top}x + c = 0 \\
    \iff ~~ & \textcolor{purple}{(T^{-1}x)}^{\top} (T^{\top}AT)\textcolor{purple}{(T^{-1}x)} + (T^{\top}b)^{\top}\textcolor{purple}{(T^{-1}x)} + c = 0 & \text{ by plain computation } \\
    \iff ~~ & y^{\top}\Lambda y + (T^{\top}b)^{\top}y + c = 0 & \text{ where } \textcolor{purple}{y = T^{-1}x} \\
    \iff ~~ & \left(\lambda_1y_1^2 + \dots + \lambda_ny_n^2\right) + \left(b_1'y_1 + \dots + b_n'y_n\right) + c = 0
\end{align*}
The point is that \textcolor{purple}{$y = T^{-1}x$} represents an orthonormal change of basis because $T \in O_n(\mathbb{R})$.  Therefore,

\begin{theorem}{Quadric Classification}
    All quadric surfaces, up to an orthonormal change of basis, look like:
    \[ \left\{ y \in \mathbb{R}^n : \left(\sum_{i = 1}^n \lambda_i y_i^2\right) + \left(\sum_{i = 1}^n b_i y_i \right) + c = 0 \right \} \] \vspace{-0.3cm}
\end{theorem}

Furthermore, if $\det A \neq 0$, then each $\lambda_i$ is nonzero, so we can complete the square, leaving $\sum_{i = 1}^n \lambda_i y_i^2 = k$ for some constant $k$.  When $k \neq 0$, we can also scale, yielding\dots

\begin{theorem}{Non-Singular Quadric Classification}
    All nonsingular quadric surfaces---those with $\det A \neq 0$ and $k \neq 0$---are affinely equivalent to:
    \[ \left\{ y \in \mathbb{R}^n : \sum_{i = 1}^n \lambda_i y_i^2 = 1 \right \}, \text{ where each $\lambda_i$ is either $+1$ or $-1$.} \] \vspace{-0.5cm}
\end{theorem}

Some corollaries of these observations:
\begin{itemize}
    \item There are only $n + 1$ non-singular quadric surfaces in $\mathbb{R}^n$.
    \item Parabolas are a degeneracy!  They correspond to a singular choice of $A$.
    \item Almost all randomly chosen $\mathbb{R}^3$ quadric surfaces (say, $67xy + 41z^2 - 14yz + 34x = 1$, for example) are affinely equivalent to one of the following: ~ $x^2 + y^2 + z^2 = 1 ~~ \text{ or } ~~ x^2 + y^2 - z^2 = 1 ~~ \text{ or } ~~ x^2 - y^2 - z^2 = 1$.
\end{itemize}




\subsection{Spectral Application \#2: Principal Component Analysis}

For this section, here's the setup:
\begin{center}
    \textbf{Setup.} Take $n$ random variables $\{X_i\}_{i = 1}^n$ in a column vector $X \in \mathbb{R}^n$. \\
    Normalize the $X_i$ so that $\mathbb{E}[X_i] = 0$ for all $i$.  (In other words, $\mathbb{E}[X] = 0$.)
\end{center}
We'll first build all necessary mathematical foundations, up to the Rayleigh-Ritz theorem.

\textbf{Definition (Variance, Covariance, Covariance Matrix).} Let $A$ and $B$ be random variables in $X$.
\begin{itemize}
    \item The \textbf{variance} of $A$ is $\Var(A) := \mathbb{E}[A^2] - \mathbb{E}[A]^2$.  \\ When $A$ is normalized so that $\mathbb{E}[A] = 0$, we have $\Var(A) = \mathbb{E}[A^2]$.
    \item The \textbf{covariance} of $A$ and $B$ is $\mathrm{Cov}(A, B) := \mathbb{E}[AB] - \mathbb{E}[A]\mathbb{E}[B]$.  \\ When $A$ and $B$ are normalized so that $\mathbb{E}[A] = \mathbb{E}[B] = 0$, we have $\mathrm{Cov}(A, B) = \mathbb{E}[AB]$.
    \item The \textbf{covariance matrix} of $X$ is $\mathrm{Cov}(X) := (\mathrm{Cov}(X_i, X_j))_{i, j}$. \\
    When $X$ is normalized so that $\mathbb{E}[X] = 0$, we have $\mathrm{Cov}(X) = \mathbb{E}[XX^{\top}]$.
\end{itemize}

\begin{theorem}{Variance of Linear Combinations}
    Let $M := \mathrm{Cov}(X)$.  Then for any $c \in \mathbb{R}^n$, we have $\Var(c^{\top}X) = c^{\top}Mc$.
\end{theorem}

\begin{proof}
    Just plain computation. \vspace{-0.1cm}
    \[ \Var(c^{\top}X) = \Var\left(\sum_{i = 1}^n c_iX_i\right) = \mathbb{E}\left [ \left(\sum_{i = 1}^n c_i X_i \right)^2 \right ] = c^{\top}Mc. ~~~ \blacksquare \] \vspace{-0.5cm}
\end{proof}

\begin{theorem}{Rayleigh-Ritz}
    Let $M$ be any symmetric matrix with normalized eigenvectors $\{v_1, \dots, v_n\}$ and corresponding eigenvalues $\lambda_1 \geq \dots \geq \lambda_n$.  Then:
    \[ \max_{c \in \mathbb{R}^n \setminus \{0\}} \dfrac{c^{\top}Mc}{c^{\top}c} = \lambda_1, ~ \text{ and more generally, } ~ \max_{\substack{c \in \mathbb{R}^n \setminus \{0\} \\ c \perp v_i \text{ for all } i < k}} \dfrac{c^{\top}Mc}{c^{\top}c} = \lambda_k. \] \vspace{-0.5cm}
\end{theorem}

\begin{proof}
    By the Spectral theorem, $\{v_1, \dots, v_n\}$ is an orthonormal basis.  Writing $c = \sum_{i = 1}^n \alpha_i v_i$ yields:
    \[ \dfrac{c^{\top}Mc}{c^{\top}c} = \dfrac{\sum_{i = 1}^n \alpha_i^2 \lambda_i}{\sum_{i = 1}^n \alpha_i^2} = \sum_{i = 1}^n \left(\dfrac{\alpha_i^2}{\sum_{j = 1}^n \alpha_j^2}\right) \lambda_i. \]
    This is a weighted average of $\{\lambda_i\}_{i = 1}^n$, and the maximum possible weighted average is $\max\{\lambda_i\}_{ i = 1}^n = \lambda_1$.
    
    More generally, when $c \perp v_i$ for all $i < k$, we're forced to have $\alpha_i = 0$ for all $i < k$.  Then we instead have a weighted average of $\{\lambda_i\}_{i = k}^n$, for which the maximum possible value is $\lambda_k$. ~ $\blacksquare$
\end{proof}

Mathematics aside, here's what we wish to do with our random variables $\{X_i\}_{i = 1}^n$ from earlier.

\begin{center}
    \textbf{Intuition.} Many samples from $X$ form a point cloud in $\mathbb{R}^n$. We wish to look at \\ this point cloud from a certain angle---a certain orthonormal basis, perhaps---so that \\ the axes of our orthonormal basis point in ``meaningful directions''.
\end{center}

What qualifies as a ``meaningful direction'', then?  The Rayleigh-Ritz theorem proposes a good candidate.

\begin{theorem}{Rayleigh-Ritz Corollary}
    Let $M$ be the covariance matrix of $X$.  Then $\frac{\Var(c^{\top}X)}{c^{\top}c}$ is maximized when $c$ is the eigenvector of $M$ with greatest eigenvalue.
\end{theorem}

Think of our point cloud as a messy ellipsoid centered at the origin.  Then the $c$ that maximizes $\frac{\Var(c^{\top}X)}{c^{\top} c}$ will be the $c$ that points along the highest-variance axis of $X$---that is, the longest axis of the ellipsoid!

In fact, the more general case of Rayleigh-Ritz indicates that the set of all eigenvectors of $M$ will point along orthogonal axes of this ellipsoid!  In conclusion, \vspace{0.1cm}

\begin{boxed-text}[32em]
    \textbf{--- Conclusion (PCA) ---} \vskip{0.5em}

    Given a point cloud $X$ of data with covariance matrix $M$, \\ we can identify the most ``principal components'' of the point cloud \\ by picking out the eigenvectors of $M$ with the greatest eigenvalues.
\end{boxed-text} \vspace{0.1cm}

Naturally, the strategy of using the eigenvectors of $M$ to analyze $X$ is called \textbf{Principal Component Analysis (PCA)}. \vspace{-0.5cm}

\begin{center}
   \frame{\includegraphics[width=0.25\textwidth]{lecture-19/PCA}}
\end{center}

In the particular context of facial recognition, \textbf{eigenfaces} are a set of eigenvectors whose linear combination with unique coefficients yields distinct faces.

\begin{center}
    \frame{\includegraphics[width=0.4\textwidth]{lecture-19/eigenfaces}}
\end{center}

\subsection{Spectral Application \#3: Singular Value Decomposition}

Here's the setup this time.
\begin{center}
    \textbf{Setup.} We have $n$ normalized random variables $\{X_i\}_{i = 1}^n$, and we sample each one $m$ times. \\
    This gives us a matrix $X \in \mathbb{R}^{m \times n}$, where $X_{i, j}$ gives the $i^{\text{th}}$ sample of $X_j$.
\end{center}

We'll first state the necessary mathematics, then show its applications, and finally prove the mathematics.

\textbf{Definition (Sample Covariance).} The \textbf{sample covariance matrix} of $X$ is the $n \times n$ matrix $\mathrm{Cov}(X) := \frac{1}{m - 1} X^{\top}X$.

Intuitively, $\mathrm{Cov}(X)_{i, j}$ gives the sample covariance of $X_i$ and $X_j$.

\begin{theorem}{Singular Value Decomposition}
    Every matrix $A \in \mathbb{R}^{m \times n}$ factors in the following way:
    \[ A = V \Sigma W^{\top} ~~ \text{ for some } ~~ V \in O_m(\mathbb{R}), ~ W \in O_n(\mathbb{R}), \text{ and diagonal }  \Sigma \in \mathbb{R}^{m \times n}. \]
    If $\Sigma$ has diagonal entries $\sigma_1 \geq \sigma_2 \geq \dots \geq 0$, then we say $\{\sigma_i\}$ are the singular values of $A$.
\end{theorem}

If $X$ factors in this way as $X = V \Sigma W^{\top}$ with singular values $\{\sigma_i\}$, then the sample covariance matrix is:
\[ \mathrm{Cov}(X) = \dfrac{1}{m - 1} (W \Sigma^{\top} V^{\top}) (V \Sigma W^{\top}) = \dfrac{1}{m - 1} W (\Sigma^{\top} \Sigma)W^{\top} = \dfrac{1}{m - 1} W \begin{bsmallmatrix} \sigma_1^2 & \cdots & 0 \\ \vdots & \ddots & \vdots \\ 0 & \cdots & \sigma_n^2 \end{bsmallmatrix}W^{\top} \]
The way $\Sigma^{\top}\Sigma$ is written above assumes $m \geq n$.  If $m < n$, then the remaining $n - m$ diagonal entries are zeroes. \vspace{0.1cm}

\begin{boxed-text}
    \textbf{--- Conclusion (SVD on Covariance Matrices) ---} \vskip{0.5em}

    Consider a matrix $X \in \mathbb{R}^{m \times n}$ of data.  Assume $m \geq n$. \\ Factor $X$ as $V \Sigma W^{\top}$ with singular values $\sigma_1 \geq \sigma_2 \geq \dots \geq 0$. \vskip{0.5em}

    Then $\mathrm{Cov}(X)$ has eigenvalues $\left\{\tfrac{\sigma_i^2}{m - 1}\right\}_{i = 1}^n$ and eigenvectors being the columns of $W$. \vskip{0.5em}

    (If $m < n$, then the remaining $n - m$ eigenvalues are zeroes.)
\end{boxed-text} \vspace{0.1cm}

Some useful diagrams depicting object-level examples of what SVD does.
\begin{center}
    \frame{\includegraphics[width=0.7\textwidth]{lecture-19/SVD}}
\end{center} \vspace{-0.3cm}
Some good things to keep in mind while we prove SVD.
\begin{itemize}
    \item For orthogonal matrices $V$ and $W$, we always have $V^{\top} = V^{-1}$ and $W^{\top} = W^{-1}$.
    \item Recall $A \in \mathbb{R}^{m \times n}$ is a map $A: \mathbb{R}^n \to \mathbb{R}^m$, so $W$ and $V$ are operators on the domain and codomain, respectively.
\end{itemize}

\begin{proof}
    Assume without loss of generality that $n \leq m$.  (If instead $n \geq m$, work with $A^{\top}$ instead of $A$.) 
    
    We begin by rephrasing our win condition.
    \begin{quote}
        \textbf{Claim.} It suffices to find orthonormal bases $\{v_1, \dots, v_m\} \subseteq \mathbb{R}^m$ and $\{w_1, \dots, w_n\} \subseteq \mathbb{R}^n$ and constants $\{\sigma_1, \dots, \sigma_n\}$ such that $Aw_i = \sigma_i v_i$ for all $i = 1, 2, \dots, n$.

        \emph{Proof:} Use the $\{v_i\}_{i = 1}^m$ and $\{w_i\}_{i = 1}^n$ as columns of $V$ and $W$, and use the $\{\sigma_i\}_{i = 1}^n$ as the diagonal of $\Sigma$.

        To show $A = V\Sigma W^{\top}$, it suffices to check that $A$ and $V \Sigma W^{\top}$ agree on the images of every element in some basis of the domain $\mathbb{R}^n$.  Let's use the basis $\{w_i\}_{i = 1}^n$.
        \[ (1) ~~~ A(w_i) = \sigma_i v_i \text{ by assumption.} ~~~~~~~~~~~ (2) ~~~ V \Sigma W^{\top}(w_i) = \sigma_i v_i \text{ by computation.} \]
        Hooray---they match. ~ $\square$
    \end{quote}

    Observe the following properties of $A^{\top}A \in \mathbb{R}^{n \times n}$.
    \begin{itemize}
        \item Since $A^{\top}A$ is symmetric, the Spectral theorem says its eigenvectors form an orthonormal basis. \\ Thus, we may take $\{w_i\}_{i = 1}^n$ to be this orthonormal basis.  \emph{(This move is forced---why?)}
        \item Since $A^{\top}A$ is positive semidefinite, the eigenvalues $\{\lambda_i\}_{i = 1}^n$ of $A^{\top}A$ are nonnegative. \\ Thus, we may take $\{\sigma_i\}_{i = 1}^n$ to be $\sigma_i = \sqrt{\lambda_i}$ for all $i$.  \emph{(This move is forced---why?)}
    \end{itemize}
    Great!  We've already found $\{w_i\}_{i = 1}^n$ and $\{\sigma_i\}_{i = 1}^n$.
    
    For $\{v_i\}_{i = 1}^m$, we'll \textcolor{orange}{just take $v_i = \frac{Aw_i}{\sigma_i}$} for $i \leq n$, then fill in the rest arbitrarily (maintaining orthonormality).

    Well, it's pretty obvious that $Aw_i = \sigma_i v_i$ will hold for all $i$.  Now we just need to show $\{v_i\}_{i = 1}^m$ is orthonormal.
    \[ \langle v_i, v_j \rangle = v_i^{\top}v_j = \frac{w_j^{\top}A^{\top}Aw_i}{\sigma_i \sigma_j} = \frac{w_j^{\top}(\sigma_i^2 w_i)}{\sigma_i \sigma_j} = \frac{\sigma_i}{\sigma_j}\langle w_j, w_i \rangle = \begin{cases} 1 & \text{ if } i = j \\ 0 & \text{ if } i \neq j \end{cases} \]
    We've checked everything we needed, so we're done. ~ \textcolor{orange}{$\clubsuit$}

   Well, okay, fine---there's one issue we missed, shown in \textcolor{orange}{orange.}  If $\sigma_i = 0$, then $v_i = \frac{Aw_i}{\sigma_i}$ doesn't make sense.  Fortunately, if $\sigma_i = 0$, there's an easy fix.
   \[ \sigma_i = 0 . ~ \implies ~ A^{\top}Aw_i = 0. ~ \implies ~ \langle w_i, A^{\top}Aw_i \rangle = 0. ~ \implies ~ \langle Aw_i, Aw_i \rangle = 0. ~ \implies ~ Aw_i = 0. \]
   Thus, the constraint $Aw_i = \sigma_i v_i$ holds trivially---again just pick $v_i$ arbitrarily, maintaining orthonormality. ~ $\blacksquare$
\end{proof}

Unlike some of our previous ``factorization'' theorems, this one applies to \emph{all} matrices!

\begin{remark}
    Despite how long the above proof is, it's actually extremely motivated---in particular, $AA^{\top}$ and $A^{\top}A$ tell us what $V$, $W$, and $\Sigma$ \emph{have} to be, and then it remains to check that everything remains consistent.
\end{remark}

\emph{Hint: It is rumored that the SVD may be useful for Problem Set \#10.}

% Then this implies we have \textbf{SVD (Singular Value Decomposition)} of our $X$, where $v_1, \dots, v_m$ and $w_1, \dots, w_n$ are orthonormal bases.
% \[ Aw_i = \sigma_i v_i \text{ and } A^{\top} v_i = \sigma_i w_i. \]
% I think the point is that it's actually very nonobvious that this implication holds.  Like i have no idea where this came from.  

% \textcolor{red}{cooked cooked cooked cooked cooked}

% We would like to show that there exist orthonormal basis $v_1, \dots, v_m$ and $w_1, \dots, w_n$ such that the above two equations hold true.

% For nonzero singular avlues, there's only one possible choice (unless you have repeated eigenvaleus???) and when the singular avlues are zero, it all works out by yourself no matter what you chose.

% what



\subsection{Linear Groups}

We're about to transition to a new topic: embedding groups into geometric objects.

\textbf{Example.} There exists a natural bijection $SO_2(\mathbb{R}) \leftrightarrow S^1$, where $S^1$ is the unit circle.

The interesting point, however, is that it's surprisingly difficult to embed $SO_3(\mathbb{R})$ into a geometric object!

Here are some failed attempts:
\begin{itemize}
    \item Every rotation in $SO_3(\mathbb{R})$ can be described by an axis $v \in S^2$ and an angle $\theta \in S^1$.  Thus, $SO_3(\mathbb{R}) \cong S^2 \times S^1$.
    \begin{itemize}
        \item Unfortunately, this ``bijection'' has a lot of messy non-injectivity!  For example, $v$ and $-v$ represent the same axis.
    \end{itemize}
    \item Every rotation in $SO_3(\mathbb{R})$ can be described using roll, pitch, and yaw.  Thus, $SO_3(\mathbb{R}) \cong S^1 \times S^1 \times S^1$.
    \begin{itemize}
        \item When built in real life, this parametrization can cause \href{https://en.wikipedia.org/wiki/Gimbal_lock}{gimbal lock} to occur!
        \item Famously, gimbal lock caused problems on the Apollo 11 Moon mission.
    \end{itemize}
\end{itemize}
Spoilers: it turns out $S^3$ and the quaternions will be the most natural parametrization of $SO_3(\mathbb{R})$.

\end{document}